\subsection{\texorpdfstring{第~\ref{sec:tangent-space}~节}{切空间}}\label{subsec:solutions-tangent-space}

\begin{description}[font=\normalfont,leftmargin=2em,labelwidth=1.5em]
\RGsolution{ex:independence-of-coordinate-tangent-vectors}{independence-of-coordinate-tangent-vectors}
考虑 $p$ 附近的参数化 $\varphi:U\subset\mathbb{R}^n\to M$ 和正文中定义的算子 $\left(\frac{\partial}{\partial x^i}\right)_p$. 若对某些 $\alpha^1,\ldots,\alpha^n\in\mathbb{R}$ 有 $\sum_{i=1}^n\alpha^i\left(\frac{\partial}{\partial x^i}\right)_p=0$, 则
\[
  \left(\sum_{i=1}^n\alpha^i\left(\frac{\partial}{\partial x^i}\right)_p\right)(f)=0
\]
对每个在 $p$ 处可微的函数 $f:M\to\mathbb{R}$ 成立. 特别地, 取 $\varphi^{-1}$ 的坐标函数, 即 $f_j:=(\varphi^{-1})^j:\varphi(U)\to\mathbb{R}$, 则 $\widehat{f_j}(x^1,\ldots,x^n)=x^j$, 从而
\[
  0=\left(\sum_{i=1}^n\alpha^i\left(\frac{\partial}{\partial x^i}\right)_p\right)(f_j)
  =\sum_{i=1}^n\alpha^i\left(\frac{\partial x^j}{\partial x^i}\right)(\varphi^{-1}(p))
  =\alpha^j,
\]
故所有 $\alpha^j$ 都等于 0.

\RGsolution{ex:chain-rule-for-manifold-maps}{chain-rule-for-manifold-maps}
设 $v\in T_pM$ 由 $v=\dot c(0)$ 给出, 其中 $c:(-\varepsilon,\varepsilon)\to M$ 是一条曲线. 则
\begin{align*}
  (\mathrm{d}(g\circ f))_p(v)
  &=\left.\frac{\mathrm{d}}{\mathrm{d}t}\right|_{t=0}(g\circ f)(c(t))\\
  &=\left.\frac{\mathrm{d}}{\mathrm{d}t}\right|_{t=0}g((f\circ c)(t))
  =(\mathrm{d}g)_{f(p)}(w),
\end{align*}
其中 $w\in T_{f(p)}N$ 是可微曲线 $(f\circ c):(-\varepsilon,\varepsilon)\to N$ 在 $t=0$ 处的切向量, 即
\[
  w=\left.\frac{\mathrm{d}}{\mathrm{d}t}\right|_{t=0}f(c(t))=(\mathrm{d}f)_p(v).
\]
因此, 对所有 $v\in T_pM$,
\[
  (\mathrm{d}(g\circ f))_p(v)=(\mathrm{d}g)_{f(p)}((\mathrm{d}f)_p(v)).
\]

\RGsolution{ex:derivative-of-antipodal-map-at-north-pole}{derivative-of-antipodal-map-at-north-pole}
将 $T_NS^n$ 和 $T_SS^n$ 与 $\mathbb{R}^{n+1}$ 中由 $x^{n+1}=0$ 给出的子空间等同, 得
\[
  (\mathrm{d}f)_N(v)
  =\left.\frac{\mathrm{d}}{\mathrm{d}t}\right|_{t=0}f(c(t))
  =\left.\frac{\mathrm{d}}{\mathrm{d}t}\right|_{t=0}(-c(t))=-\dot c(0)=-v
\]
(其中 $c:(-\varepsilon,\varepsilon)\to S^n$ 满足 $\dot c(0)=v$), 即 $(\mathrm{d}f)_N$ 就是乘以 $-1$. 另一种做法是, 使用 $N$ 和 $S$ 处的参数化 $(\mathbb{R}^n,\pi_S^{-1})$ 和 $(\mathbb{R}^n,\pi_N^{-1})$, 有
\begin{align*}
  \widehat f(y)&=\pi_N\circ f\circ\pi_S^{-1}(y)
  =\pi_N\circ f\left(\frac{2y}{1+\lVert y\rVert^2},\frac{1-\lVert y\rVert^2}{1+\lVert y\rVert^2}\right)\\
  &=\pi_N\left(-\frac{2y}{1+\lVert y\rVert^2},\frac{\lVert y\rVert^2-1}{1+\lVert y\rVert^2}\right)=-y.
\end{align*}
这个映射在 $\pi_S(N)=0$ 处的 Jacobian 矩阵是 $(\mathrm{d}\widehat f)_0=-\operatorname{id}$. 因此, 若 $y$ 是第一个参数化对应的局部坐标, 而 $z$ 是第二个参数化对应的局部坐标, 则
\[
  (\mathrm{d}f)_N\left(v^1\frac{\partial}{\partial y^1}+\cdots+v^n\frac{\partial}{\partial y^n}\right)
  =-v^1\frac{\partial}{\partial z^1}-\cdots-v^n\frac{\partial}{\partial z^n}.
\]

\RGsolution{ex:coordinate-differential-of-scalar-function}{coordinate-differential-of-scalar-function}
设 $c:(-\varepsilon,\varepsilon)\to W$ 是 $W\subset M$ 中一条满足 $c(0)=p$ 的曲线, 并令 $v:=\dot c(0)$. 则在坐标图 $x:W\to\mathbb{R}^n$ 上,
\[
  (\mathrm{d}f)_pv=\frac{\mathrm{d}(f\circ c)}{\mathrm{d}t}(0)
  =\sum_{i=1}^n\dot{x}^i(0)\frac{\partial\widehat f}{\partial x^i}(x(p)),
\]
其中在局部坐标下写作 $\widehat c(t)=(x^1(t),\ldots,x^n(t))$. 另一方面,
\[
  (\mathrm{d}x^j)_pv=\frac{\mathrm{d}(x^j\circ c)}{\mathrm{d}t}(0)
  =\sum_{i=1}^n\dot{x}^i(0)\frac{\partial x^j}{\partial x^i}(x(p))=\dot{x}^j(0),
\]
所求结论随即成立.

\RGsolution{ex:tangent-bundle-chart-construction}{tangent-bundle-chart-construction}
显然,
\[
  \bigcup_\alpha\Phi_\alpha(U_\alpha\times\mathbb{R}^n)
  =\bigcup_{p\in\bigcup_\alpha\varphi_\alpha(U_\alpha)}T_pM=TM.
\]
此外, 容易看出, 这些参数化在 $TM$ 上定义的拓扑满足 Hausdorff 公理和第二可数公理. 最后, 若 $W=\Phi_\alpha(U_\alpha)\cap\Phi_\beta(U_\beta)\ne\varnothing$, 则过渡映射
\[
  \Phi_\beta^{-1}\circ\Phi_\alpha:\Phi_\alpha^{-1}(W)\to\Phi_\beta^{-1}(W)
\]
和
\[
  \Phi_\alpha^{-1}\circ\Phi_\beta:\Phi_\beta^{-1}(W)\to\Phi_\alpha^{-1}(W)
\]
分别由
\[
  (\Phi_\beta^{-1}\circ\Phi_\alpha)(x,v)
  =\bigl((\varphi_\beta^{-1}\circ\varphi_\alpha)(x),(\mathrm{d}(\varphi_\beta^{-1}\circ\varphi_\alpha))_x(v)\bigr)
\]
和
\[
  (\Phi_\alpha^{-1}\circ\Phi_\beta)(x,v)
  =\bigl((\varphi_\alpha^{-1}\circ\varphi_\beta)(x),(\mathrm{d}(\varphi_\alpha^{-1}\circ\varphi_\beta))_x(v)\bigr)
\]
给出, 因而可微 (上述公式使用了标准等同 $T_x\mathbb{R}^n\cong\mathbb{R}^n$).
\end{description}
